\section{The backdoored models}
\label{app:panel}

The 65 models behind every detection number in the paper are those in \cref{tab:panel} whose attack success rate is at least 0.85. Models below that bar are marked with a dagger and take no part in the results, since a detector cannot be credited or blamed for a backdoor that was never planted. Each row is 1 dataset and 1 attack, and the 3 columns are the poison rates 1\%, 5\% and 10\%. A model that clears the bar prints its attack success rate over its clean accuracy. A model below the bar prints only its attack success rate.

Reading down the table shows what the evaluation can and cannot say. BadNets, Blend, BPP and LF implant at every rate on every dataset. TaCT implants everywhere except GTSRB at 1\%. WaNet implants only at 5\% and 10\% and not on CIFAR-100, SIG only on CIFAR-10 at 10\%, and Adaptive-Blend never reaches the bar. The hard-attack subset of the main text is therefore BPP, WaNet, TaCT and SIG. \Cref{tab:panel-benign} gives the clean accuracy of the benign reference model of each dataset, trained with the same recipe, which every clean-accuracy drop in the paper subtracts against.

% generated by scripts/paper/tab_panel.py from results/coverage/coverage.json at 2026-09-23T14:42:18.773755+00:00, commit 039fc6fea9212f8f6d49218ed14b336aa8f6a0d3-dirty. Do not edit by hand.
\begin{table}[htbp]
\centering
\small
\caption{The ViT panel, from results/coverage/coverage.json: 71 of 105 declared models clear the attack success bar 0.85 and carry every table in this paper. A model that clears prints attack success rate over clean accuracy. A model below the bar prints only its attack success rate, marked $^\dagger$, so its exclusion stays visible.}
\label{tab:panel}
\begin{adjustbox}{max width=\linewidth}
\begin{tabular}{llrrr}
\toprule
dataset & attack & 1\% & 5\% & 10\% \\
\midrule
CIFAR-10 & BadNets & 0.997 / 0.952 & 1.000 / 0.945 & 1.000 / 0.942 \\
CIFAR-10 & Blend & 1.000 / 0.954 & 1.000 / 0.957 & 1.000 / 0.951 \\
CIFAR-10 & LF & 0.982 / 0.953 & 0.997 / 0.944 & 0.999 / 0.951 \\
CIFAR-10 & BPP & 0.982 / 0.952 & 0.987 / 0.950 & 0.994 / 0.954 \\
CIFAR-10 & WaNet & 0.112$^\dagger$ & 0.961 / 0.913 & 0.890 / 0.946 \\
CIFAR-10 & TaCT & 0.985 / 0.938 & 0.996 / 0.946 & 1.000 / 0.858 \\
CIFAR-10 & SIG & 0.340$^\dagger$ & 0.599$^\dagger$ & 0.901 / 0.846 \\
CIFAR-10 & Adaptive-Blend & 0.622$^\dagger$ & 0.594$^\dagger$ & 0.622$^\dagger$ \\
CIFAR-100 & BadNets & 1.000 / 0.825 & 1.000 / 0.825 & 1.000 / 0.821 \\
CIFAR-100 & Blend & 0.989 / 0.829 & 1.000 / 0.830 & 1.000 / 0.816 \\
CIFAR-100 & LF & 0.942 / 0.830 & 0.953 / 0.829 & 0.995 / 0.823 \\
CIFAR-100 & BPP & 0.914 / 0.824 & 0.988 / 0.833 & 0.991 / 0.817 \\
CIFAR-100 & WaNet & 0.057$^\dagger$ & 0.643$^\dagger$ & 0.793$^\dagger$ \\
CIFAR-100 & TaCT & 0.989 / 0.820 & 1.000 / 0.825 & 0.989 / 0.821 \\
CIFAR-100 & SIG & 0.250$^\dagger$ & 0.249$^\dagger$ & 0.178$^\dagger$ \\
CIFAR-100 & Adaptive-Blend & 0.536$^\dagger$ & 0.579$^\dagger$ & 0.604$^\dagger$ \\
GTSRB & BadNets & 1.000 / 0.994 & 0.573$^\dagger$ & 1.000 / 0.972 \\
GTSRB & Blend & 1.000 / 0.991 & 0.951$^\dagger$ & 1.000 / 0.988 \\
GTSRB & LF & 0.979 / 0.990 & 0.999 / 0.974 & 0.999 / 0.988 \\
GTSRB & BPP & 0.868 / 0.986 & 0.991 / 0.978 & 0.995 / 0.986 \\
GTSRB & WaNet & 0.119$^\dagger$ & 0.830$^\dagger$ & 0.947 / 0.958 \\
GTSRB & TaCT & 0.000$^\dagger$ & 1.000 / 0.986 & 1.000 / 0.931 \\
GTSRB & SIG & 0.356$^\dagger$ & 0.673$^\dagger$ & -- \\
GTSRB & Adaptive-Blend & 0.560$^\dagger$ & 0.777$^\dagger$ & 0.838$^\dagger$ \\
Tiny ImageNet & BadNets & 0.999 / 0.752 & 1.000 / 0.752 & 1.000 / 0.752 \\
Tiny ImageNet & Blend & 0.999 / 0.760 & 0.999 / 0.754 & 1.000 / 0.749 \\
Tiny ImageNet & LF & 0.922 / 0.762 & 0.987 / 0.753 & 0.973 / 0.740 \\
Tiny ImageNet & BPP & 0.971 / 0.748 & 0.985 / 0.746 & 0.996 / 0.757 \\
Tiny ImageNet & WaNet & 0.178$^\dagger$ & 0.922 / 0.748 & 0.968 / 0.747 \\
Tiny ImageNet & TaCT & 0.976 / 0.750 & 0.952 / 0.743 & 0.929 / 0.756 \\
Tiny ImageNet & SIG & 0.076$^\dagger$ & 0.432$^\dagger$ & 0.122$^\dagger$ \\
Tiny ImageNet & Adaptive-Blend & 0.554$^\dagger$ & 0.783$^\dagger$ & 0.505$^\dagger$ \\
SVHN & BadNets & -- & -- & -- \\
SVHN & Blend & -- & -- & 1.000 / 0.951 \\
SVHN & LF & -- & -- & -- \\
SVHN & BPP & -- & -- & -- \\
SVHN & WaNet & -- & -- & -- \\
SVHN & TaCT & -- & -- & -- \\
SVHN & SIG & 0.700$^\dagger$ & 0.880 / 0.958 & 0.971 / 0.958 \\
SVHN & Adaptive-Blend & -- & -- & -- \\
EuroSAT & BadNets & -- & -- & -- \\
EuroSAT & Blend & -- & -- & 1.000 / 0.980 \\
EuroSAT & LF & -- & -- & -- \\
EuroSAT & BPP & -- & -- & -- \\
EuroSAT & WaNet & -- & -- & -- \\
EuroSAT & TaCT & -- & -- & -- \\
EuroSAT & SIG & 0.813$^\dagger$ & 0.873 / 0.979 & 0.922 / 0.972 \\
EuroSAT & Adaptive-Blend & -- & -- & -- \\
\bottomrule
\end{tabular}
\end{adjustbox}
\end{table}

% generated by scripts/paper/tab_panel.py from results/coverage/coverage.json at 2026-09-11T04:15:50.733381+00:00, commit 4a04e5dd204b14d8368fb7eeef7092bd331e00d9-dirty. Do not edit by hand.
\begin{table}[htbp]
\centering
\small
\caption{Benign-model clean accuracy per dataset, the reference every clean-accuracy-drop figure in this paper subtracts against. Same architecture, optimizer and 15 epochs as every backdoored model in the panel.}
\label{tab:panel-benign}
\begin{adjustbox}{max width=\linewidth}
\begin{tabular}{lr}
\toprule
dataset & benign clean accuracy \\
\midrule
CIFAR-10 & 0.953 \\
CIFAR-100 & 0.811 \\
GTSRB & 0.991 \\
Tiny ImageNet & 0.755 \\
SVHN & 0.964 \\
EuroSAT & 0.977 \\
\bottomrule
\end{tabular}
\end{adjustbox}
\end{table}

\FloatBarrier

\section{The placement basis}
\label{app:basis}

\Cref{tab:basis-declaration} lists the 18 placements every model was swept over. A placement is a site (where the perturbation is injected, \cref{fig:vit-block}) paired with a perturbation operator (what the perturbation does), optionally restricted to a band of blocks. The family column groups sites by where they sit relative to the residual stream. Input-side sites perturb a sublayer's input before its LayerNorm. Residual-adjacent sites perturb a branch output just before its residual add or the stream just after it. The structured sites perturb attention heads and MLP hidden units, and the ported sites come from other detectors (scaling input pixels as SCALE-UP does, scaling a LayerNorm's gain as IBD-PSC does). The rates column is the ladder each placement was swept over. Masks and dropout use a share of the tensor, Gaussian noise a multiple of the tensor's standard deviation.

% generated by scripts/paper/app_basis.py from configs/psbd_basis.json at 2026-09-10T23:32:03.398466+00:00, commit 831df7ef5e8bc9285264b261f848e1f2cad5ef61-dirty. Do not edit by hand.
\begin{table}[t]
\centering
\small
\caption{The placement basis as configs/psbd\_basis.json declares it: every panel cell carries every 1 of these placements, each swept over the rate ladder in the last column.}
\label{tab:basis-declaration}
\begin{tabular}{llllll}
\toprule
placement id & position & operator & blocks & family & rates \\
\midrule
before\_attention\_norm\_token\_mask & before\_attention\_norm & token\_mask & all & input\_side & 10 (0.05 to 0.9) \\
before\_attention\_norm\_channel\_mask & before\_attention\_norm & channel\_mask & all & input\_side & 10 (0.05 to 0.9) \\
before\_attention\_norm\_gaussian & before\_attention\_norm & gaussian & all & input\_side & 12 (0.05 to 8) \\
before\_mlp\_norm\_token\_mask & before\_mlp\_norm & token\_mask & all & input\_side & 10 (0.05 to 0.9) \\
both\_sublayer\_inputs\_token\_mask & both\_sublayer\_inputs & token\_mask & all & input\_side & 9 (0.1 to 0.9) \\
input\_pixels\_scale\_up & input\_pixels & scale\_up & all & input\_side & 11 (0.5 to 25) \\
before\_mlp\_gaussian & before\_mlp & gaussian & all & input\_side & 12 (0.05 to 8) \\
before\_attention\_residual\_token\_mask & before\_attention\_residual & token\_mask & all & residual\_adjacent & 10 (0.05 to 0.9) \\
after\_attention\_residual\_token\_mask & after\_attention\_residual & token\_mask & all & residual\_adjacent & 10 (0.05 to 0.9) \\
pre\_residual & pre\_residual & dropout & all & residual\_adjacent & 9 (0.1 to 0.9) \\
post\_residual & post\_residual & dropout & all & residual\_adjacent & 16 (0.005 to 0.9) \\
mlp\_neurons\_channel\_mask & mlp\_neurons & channel\_mask & all & structured & 10 (0.05 to 0.9) \\
mlp\_norm\_out\_gain\_scale & mlp\_norm\_out & gain\_scale & all & ported & 11 (0.25 to 15) \\
before\_attention\_norm\_blocks\_5\_8\_token\_mask & before\_attention\_norm & token\_mask & 5-8 & depth\_band & 12 (0.05 to 0.99) \\
before\_attention\_norm\_blocks\_9\_12\_token\_mask & before\_attention\_norm & token\_mask & 9-12 & depth\_band & 12 (0.05 to 0.99) \\
pre\_residual\_blocks\_1\_4 & pre\_residual & dropout & 1-4 & depth\_band & 11 (0.1 to 0.99) \\
pre\_residual\_blocks\_5\_8 & pre\_residual & dropout & 5-8 & depth\_band & 11 (0.1 to 0.99) \\
pre\_residual\_blocks\_9\_12 & pre\_residual & dropout & 9-12 & depth\_band & 11 (0.1 to 0.99) \\
\bottomrule
\end{tabular}
\end{table}

\FloatBarrier

\Cref{tab:basis-ranking} ranks the same 18 placements on the 65 models. For each placement the table gives its mean AUROC at the matched rule and at the adaptive rule (\cref{sec:results:reading}), its lowest single-model AUROC at the adaptive rule, and how many models fall below chance. The top 3 placements all mask whole tokens on an input to attention. The gain-scaling and noise placements sit in the bottom half, and scaling input pixels is last. Attention-input token masking is second by 0.001 at the adaptive rule behind token masking on the attention branch before its residual add, and first at the matched rule. We recommend it because it leads at matched disturbance and because its site has the mechanistic reading of \cref{sec:mechanism}. \Cref{fig:basis} plots this table.

% generated by scripts/paper/app_basis.py from results/coverage/coverage.json, configs/psbd_basis.json, results/<folder>/psbd_metrics.json (71 cells) at 2026-09-23T16:00:54.049534+00:00, commit afe5fd94d50b77f30f654fe075b768040987adef-dirty. Do not edit by hand.
\begin{table}[htbp]
\centering
\small
\caption{Every placement of the basis ranked by mean AUROC on ViT-B/16. n is the number of models whose rate ladder reaches the shift ratio target, and the last column is the lowest single-model AUROC.}
\label{tab:basis-ranking}
\begin{adjustbox}{max width=\linewidth}
\begin{tabular}{rllrrr}
\toprule
Rank & Placement & Family & n & AUROC & Lowest AUROC \\
\midrule
1 & token mask, attention input & input side & 69 & 0.927 & 0.418 \\
2 & token mask, attention output before the add & residual adjacent & 68 & 0.927 & 0.537 \\
3 & token mask, attention input, blocks 1 to 4 & depth band & 12 & 0.924 & 0.653 \\
4 & token mask, both sublayer inputs & input side & 71 & 0.898 & 0.466 \\
5 & gain scale, MLP norm output & ported & 70 & 0.896 & 0.266 \\
6 & token mask, attention input, blocks 9 to 12 & depth band & 31 & 0.893 & 0.514 \\
7 & token mask, attention input, blocks 5 to 8 & depth band & 44 & 0.883 & 0.244 \\
8 & channel mask, attention input & input side & 71 & 0.883 & 0.091 \\
9 & dropout, before both residual adds, blocks 9 to 12 & depth band & 46 & 0.880 & 0.372 \\
10 & dropout, before both residual adds, blocks 5 to 8 & depth band & 71 & 0.873 & 0.365 \\
11 & noise, MLP input after norm & input side & 71 & 0.871 & 0.226 \\
12 & dropout, attention input & input side & 65 & 0.859 & 0.074 \\
13 & dropout, before both residual adds & residual adjacent & 71 & 0.838 & 0.093 \\
14 & token mask, MLP input & input side & 71 & 0.826 & 0.048 \\
15 & dropout, after both residual adds & residual adjacent & 71 & 0.818 & 0.222 \\
16 & dropout, attention input after norm & input side & 65 & 0.811 & 0.165 \\
17 & dropout, embedding output & input side & 65 & 0.810 & 0.172 \\
18 & channel mask, MLP neurons & structured & 71 & 0.808 & 0.250 \\
19 & noise, attention output before the add & residual adjacent & 65 & 0.806 & 0.055 \\
20 & dropout, before both residual adds, blocks 1 to 4 & depth band & 71 & 0.805 & 0.080 \\
21 & dropout, stream after the attention add & residual adjacent & 65 & 0.804 & 0.190 \\
22 & noise, attention input & input side & 71 & 0.804 & 0.031 \\
23 & channel mask, attention output before the add & residual adjacent & 65 & 0.789 & 0.029 \\
24 & token mask, stream after the attention add & residual adjacent & 71 & 0.761 & 0.050 \\
25 & dropout, MLP input & input side & 65 & 0.759 & 0.070 \\
26 & scale up, input pixels & input side & 71 & 0.756 & 0.225 \\
27 & dropout, attention output before the add & residual adjacent & 0 & -- & -- \\
\bottomrule
\end{tabular}
\end{adjustbox}
\end{table}

\FloatBarrier

\section{Depth bands}
\label{app:bands}

\Cref{tab:depth-bands} is the measurement behind \cref{sec:results:depth}. Attention-input token masking and dropout before the residual add were each rerun restricted to a third of the network (blocks 1 to 4, 5 to 8 and 9 to 12) and compared, at the matched rule, with their all-blocks versions on the same models. The delta column is the paired difference from the all-blocks version with its bootstrap interval. The last column is the clean shift ratio the chosen rate reached, which confirms that the comparison is at matched disturbance. Restricting token masking to any band costs about 0.06. Restricting residual dropout to the middle band gains 0.05 and to the early band loses 0.06. The early band of token masking was not in the original basis and is being swept.

% generated by scripts/paper/tab_depth_bands.py from results/coverage/coverage.json, results/<folder>/psbd_metrics.json (65 cells) at 2026-09-11T04:12:43.663538+00:00, commit 4a04e5dd204b14d8368fb7eeef7092bd331e00d9-dirty. Do not edit by hand.
\begin{table}[htbp]
\centering
\small
\caption{Depth-band placements at the matched 0.6 rate, mean AUROC at q0.25, the paired delta against the family's all-blocks placement with a 5000-resample bootstrap 95\% interval, and the mean achieved clean-validation shift ratio at the chosen rate. Token masking at the attention input restricted to blocks 1 to 4 has not been swept yet on any model in the panel of 65 backdoored models and prints pending throughout.}
\label{tab:depth-bands}
\begin{adjustbox}{max width=\linewidth}
\begin{tabular}{lrrrlr}
\toprule
placement & n & mean AUROC matched06 & delta vs all blocks & 95\% CI & mean achieved shift \\
\midrule
input-side, all blocks & 65 & 0.926 & -- & -- & 0.595 \\
input-side, blocks 5-8 & 65 & 0.863 & -0.063 & [-0.100, -0.031] & 0.603 \\
input-side, blocks 9-12 & 65 & 0.864 & -0.062 & [-0.108, -0.021] & 0.489 \\
input-side, blocks 1-4 & \pending & \pending & \pending & \pending & \pending \\
residual-adjacent, all blocks & 65 & 0.841 & -- & -- & 0.658 \\
residual-adjacent, blocks 1-4 & 65 & 0.779 & -0.062 & [-0.081, -0.042] & 0.582 \\
residual-adjacent, blocks 5-8 & 65 & 0.893 & +0.052 & [+0.027, +0.079] & 0.565 \\
residual-adjacent, blocks 9-12 & 65 & 0.874 & +0.033 & [-0.012, +0.076] & 0.598 \\
\bottomrule
\end{tabular}
\end{adjustbox}
\end{table}

\FloatBarrier

\section{Per-layer profiles of the backdoor direction}
\label{app:layers}

\Cref{tab:tac-layers} gives, for each of the 18 models in \cref{fig:tac-layers}, the block at which the backdoor direction's norm relative to the clean activation peaks, that peak value, the first block at which the relative norm reaches half its final value (the onset), the final-layer linear centered kernel alignment (CKA) between clean and triggered activations (1.0 means the trigger leaves the representation unchanged), and the mean absolute trigger-activated change (TAC) per dimension at the final layer. The benign rows are the control. Their peak relative norm never exceeds 0.12 and their CKA is 1.0, so the profiles of the backdoored rows are learned structure and not an artefact of the trigger's pixels. TaCT stands out with a peak of 0.26 to 0.38 and a CKA of 0.98 to 0.99. A source-specific attack changes the representation of the few source classes only, so a direction averaged over all classes is small even though the attack succeeds.

% generated by scripts/paper/mech_tac_layers.py from results/_experiments/tac_layers/tac_layers.json at 2026-09-23T18:21:02.095010+00:00, commit 5a1f1c5814dc5676c093f3021490a849712a13f4-dirty. Do not edit by hand.
\begin{table}[htbp]
\centering
\small
\caption{Per-layer profile of the backdoor direction on 18 ViT-B/16 models, 500 paired images each.}
\label{tab:tac-layers}
\begin{adjustbox}{max width=\linewidth}
\begin{tabular}{llrrrrr}
\toprule
Dataset & Model & Peak layer & Peak rel. norm & Onset layer & Final CKA & Final TAC \\
\midrule
GTSRB & BadNets 10\% & 12 & 1.64 & 7 & 0.57 & 1.05 \\
GTSRB & benign & 5 & 0.01 & 2 & 1.00 & 0.01 \\
GTSRB & Blend 10\% & 12 & 3.40 & 9 & 0.12 & 0.62 \\
GTSRB & BPP 10\% & 9 & 1.66 & 7 & 0.41 & 0.76 \\
GTSRB & LF 10\% & 12 & 1.68 & 7 & 0.48 & 0.86 \\
GTSRB & TaCT 10\% & 12 & 0.38 & 12 & 0.99 & 0.20 \\
GTSRB & WaNet 10\% & 12 & 1.38 & 8 & 0.29 & 0.59 \\
CIFAR-100 & BadNets 1\% & 12 & 0.79 & 8 & 0.44 & 0.57 \\
CIFAR-100 & BadNets 10\% & 10 & 0.92 & 8 & 0.39 & 0.57 \\
CIFAR-100 & benign & 10 & 0.12 & 6 & 0.99 & 0.06 \\
CIFAR-100 & Blend 1\% & 10 & 1.04 & 7 & 0.23 & 0.62 \\
CIFAR-100 & Blend 10\% & 9 & 1.70 & 5 & 0.16 & 0.71 \\
CIFAR-100 & BPP 10\% & 12 & 1.28 & 8 & 0.20 & 0.75 \\
CIFAR-100 & LF 10\% & 12 & 1.46 & 8 & 0.24 & 0.75 \\
CIFAR-100 & TaCT 10\% & 10 & 0.26 & 9 & 0.98 & 0.13 \\
Tiny ImageNet & BadNets 10\% & 12 & 0.99 & 9 & 0.41 & 0.55 \\
Tiny ImageNet & benign & 9 & 0.04 & 5 & 1.00 & 0.02 \\
Tiny ImageNet & Blend 10\% & 10 & 1.52 & 7 & 0.15 & 0.66 \\
\bottomrule
\end{tabular}
\end{adjustbox}
\end{table}

\FloatBarrier

\section{Seed replicates}
\label{app:seeds}

\Cref{tab:seeds} gives the 3-seed readings behind \cref{sec:robustness:seeds}. Each row is a backdoored model trained 3 times with different seeds. The columns give the smallest attack success rate among the 3 runs, the AUROC of attention-input token masking on each seed with its standard deviation, the same for post-residual dropout, and the mean and standard deviation of their paired gain. Token masking's standard deviation stays under 0.04 on every row, while post-residual dropout's reaches 0.28 on CIFAR-100 BadNets at 5\%. Further seeds are being trained for the hard attacks and the other datasets.

% generated by scripts/paper/tab_seeds.py from results/coverage/coverage.json, results/<folder>[_seed_N]/psbd_metrics.json (14 cells), checkpoints/<folder>_seed_N/args.json at 2026-09-23T14:45:07.243909+00:00, commit 039fc6fea9212f8f6d49218ed14b336aa8f6a0d3-dirty. Do not edit by hand.
\begin{table}[htbp]
\centering
\small
\caption{Seed replicates of PSBD-TM and PSBD-RD on the models trained at 3 seeds. Min ASR is the lowest attack success among the 3 runs.}
\label{tab:seeds}
\begin{adjustbox}{max width=\linewidth}
\begin{tabular}{lllrlrlrrr}
\toprule
Dataset & Attack & Rate & Min ASR & PSBD-TM, seeds 0 / 1 / 2 & sd & PSBD-RD, seeds 0 / 1 / 2 & sd & Gain mean & Gain sd \\
\midrule
CIFAR-100 & BadNets & 0.01 & 0.999 & 0.988 / 0.996 / 0.984 & 0.006 & 0.744 / 0.512 / 0.460 & 0.151 & +0.417 & 0.151 \\
CIFAR-100 & BadNets & 0.05 & 1.000 & 0.994 / 0.996 / 0.950 & 0.026 & 0.826 / 0.493 / 0.274 & 0.278 & +0.449 & 0.259 \\
CIFAR-100 & BadNets & 0.1 & 1.000 & 0.998 / 0.998 / 0.995 & 0.002 & 0.497 / 0.659 / 0.565 & 0.081 & +0.423 & 0.081 \\
CIFAR-100 & Blend & 0.01 & 0.988 & 0.976 / 0.991 / 0.923 & 0.036 & 0.792 / 0.991 / 0.871 & 0.100 & +0.079 & 0.095 \\
CIFAR-100 & Blend & 0.05 & 0.999 & 0.979 / 0.995 / 0.982 & 0.009 & 0.999 / 0.996 / 0.902 & 0.055 & +0.020 & 0.053 \\
CIFAR-100 & Blend & 0.1 & 0.999 & 0.998 / 0.996 / 0.989 & 0.005 & 1.000 / 0.995 / 0.992 & 0.004 & -0.001 & 0.003 \\
Tiny ImageNet & BadNets & 0.01 & 1.000 & 0.980 / 0.976 / 0.985 & 0.004 & 0.892 / 0.828 / 0.908 & 0.042 & +0.104 & 0.038 \\
Tiny ImageNet & BadNets & 0.05 & 1.000 & 0.984 / 0.992 / 0.975 & 0.008 & 0.837 / 0.906 / 0.696 & 0.107 & +0.171 & 0.099 \\
Tiny ImageNet & BadNets & 0.1 & 1.000 & 0.993 / 0.986 / 0.969 & 0.012 & 0.973 / 0.856 / 0.650 & 0.164 & +0.156 & 0.151 \\
Tiny ImageNet & Blend & 0.01 & 1.000 & 0.992 / 0.999 / 0.996 & 0.003 & 0.996 / 0.999 / 0.997 & 0.001 & -0.001 & 0.002 \\
Tiny ImageNet & Blend & 0.05 & 0.997 & 0.997 / 0.998 / 0.979 & 0.011 & 0.999 / 1.000 / 0.994 & 0.003 & -0.006 & 0.007 \\
Tiny ImageNet & Blend & 0.1 & 1.000 & 0.998 / 0.997 / 0.997 & 0.001 & 0.999 / 1.000 / 1.000 & 0.000 & -0.002 & 0.001 \\
Tiny ImageNet & WaNet & 0.05 & 0.828 & 0.930 / 0.950 / 0.874 & 0.039 & 0.953 / 0.961 / 0.877 & 0.047 & -0.013 & 0.011 \\
Tiny ImageNet & WaNet & 0.1 & 0.931 & 0.950 / 0.967 / 0.948 & 0.010 & 0.966 / 0.974 / 0.928 & 0.024 & -0.001 & 0.018 \\
\bottomrule
\end{tabular}
\end{adjustbox}
\end{table}

\FloatBarrier

\section{The other detectors on the smoke-test models}
\label{app:smoke}

\Cref{tab:detector-smoke} is the smoke test behind \cref{sec:robustness:detectors}. 10 input-level detectors, each ported and checked against its reference implementation, run on a benign, a BadNets and a Blend GTSRB model at 5\% poisoning with 500 validation images, on the same splits and threshold rule as PSBD. The passes column is the number of forward passes each detector spends per input and the last column the measured seconds per input on 1 A100. Every detector reads chance on the benign model, which is the check that its sign convention is right. On BadNets STRIP, CD-L, TED and TeCo reach 0.99 or more, so a patch trigger at 5\% is easy for every consistency detector. The ranking on the hard attacks is what the full run will decide. The 2 detectors that fail, IBD-PSC and SentiNet, fail for ViT-specific reasons stated in \cref{sec:robustness:detectors}.

% generated by scripts/paper/tab_detector_smoke.py from docs/runs/2026-09-10-detector-smoke.md at 2026-09-10T23:37:23.095504+00:00, commit 831df7ef5e8bc9285264b261f848e1f2cad5ef61-dirty. Do not edit by hand.
\begin{table}[t]
\centering
\small
\caption{Competitor detectors on the login-node smoke run: 3 GTSRB checkpoints at the middle panel rate, AUROC at the headline quantile on the benign reference and the 2 backdoored models, TPR at the smallest budget on the backdoored ones, forward passes per input and seconds per input. Acceptance only, the panel sweep is pending.}
\label{tab:detector-smoke}
\begin{tabular}{lrrrrrrr}
\toprule
detector & passes & AUROC benign & AUROC BadNet & AUROC Blend & TPR@1\% BadNet & TPR@1\% Blend & s per input \\
\midrule
confidence & 1 & 0.503 & 0.738 & 0.876 & 0.000 & 0.034 & 0.001 \\
strip & 8 & 0.500 & 1.000 & 0.976 & 1.000 & 0.905 & 0.005 \\
scale\_up & 6 & 0.499 & 0.914 & 0.689 & 0.000 & 0.000 & 0.004 \\
scale\_up\_data\_limited & 6 & 0.499 & 0.909 & 0.678 & 0.000 & 0.000 & 0.002 \\
ibd\_psc & 6 & 0.500 & 0.481 & 0.970 & 0.000 & 0.724 & 0.002 \\
teco & 71 & 0.495 & 0.994 & 0.922 & 0.757 & 0.258 & 0.057 \\
cd\_l & 251 & 0.493 & 1.000 & 0.962 & 1.000 & 0.596 & 0.112 \\
beatrix & 1 & 0.510 & 0.453 & 0.819 & 0.000 & 0.000 & 0.001 \\
ted & 1 & 0.509 & 0.999 & 0.990 & 0.990 & 0.811 & 0.001 \\
sentinet & 202 & 0.499 & 0.107 & 0.379 & 0.056 & 0.051 & 0.057 \\
\bottomrule
\end{tabular}
\end{table}

